\section{Franck--Condon 原理}
\label{ch:franck-condon}

前两章先从精确 Born--Huang 多通道理论降到单一绝热势能面，再把这一单势能面核哈密顿量特殊化为双原子分子的振动--转动问题。本章转而讨论两个不同电子势能面之间的光学跃迁。关键问题是：电子跃迁虽然很快，但核波函数并不会凭空消失；它在两个势能面上的投影决定了振动态谱线强度。单模位移谐振子只是最简单的可解模型，真实多原子分子还会出现频率改变和正规-模转动。
\subsection{从电子跃迁矩阵元到 Condon 近似}

在 BO 表示中，把初末态写成

\[
|g,\nu\rangle
=|\phi_g(\mathbf r;\mathbf Q)\rangle
|\chi_{g\nu}(\mathbf Q)\rangle,
\]

\[
|e,\nu'\rangle
=|\phi_e(\mathbf r;\mathbf Q)\rangle
|\chi_{e\nu'}(\mathbf Q)\rangle,
\]

其中 \(\mathbf Q\) 是质量加权正规坐标。电偶极矩阵元为

\[
M_{\nu'\nu}
=\int\dd\mathbf Q\,
\chi_{e\nu'}^*(\mathbf Q)
\mu_{eg}(\mathbf Q)
\chi_{g\nu}(\mathbf Q),
\]

电子跃迁偶极

\[
\mu_{eg}(\mathbf Q)
\equiv
\langle\phi_e(\mathbf Q)|\hat\mu|\phi_g(\mathbf Q)\rangle_{\mathbf r}.
\]

在参考构型 \(\mathbf Q=0\) 附近作 Taylor 展开：

\begin{equation}
\mu_{eg}(\mathbf Q)
=\mu_{eg}(0)
+\sum_\alpha
\left(\frac{\partial\mu_{eg}}{\partial Q_\alpha}\right)_0Q_\alpha
+O(Q^2).
\label{eq:herzberg-teller-expansion}
\end{equation}

Condon 近似是保留第一项：

\[
\mu_{eg}(\mathbf Q)\simeq\mu_{eg}(0).
\]

于是

\begin{equation}
M_{\nu'\nu}
\simeq
\mu_{eg}(0)
\langle\chi_{e\nu'}|\chi_{g\nu}\rangle,
\qquad
I_{\nu'\nu}\propto
|\langle\chi_{e\nu'}|\chi_{g\nu}\rangle|^2.
\label{eq:franck-condon-factor}
\end{equation}

若 \(\mu_{eg}(0)=0\) 是由对称性强制的，线性 Herzberg--Teller 项反而成为最低非零阶，此时“谱线强度只等于振动重叠平方”不再成立。因此 Franck--Condon 规则本身也有明确的截断条件。

Condon 近似把跃迁矩元简化为振动重叠问题，但重叠积分的计算仍需知道两个势能面的几何关系。最简单的情形是单模、等频率位移——此时 Franck--Condon 因子完全由 Huang--Rhys 因子 $S_{HR}$ 决定，且激发态振动真空在基态基底中恰为相干态。

\subsection{Huang--Rhys 因子}
先看一维、两条势能面曲率相同的模型。此小节的 $Q$ 是有长度量纲的物理坐标，故动能为 $P^2/(2M)$；多模计算时再统一作质量加权：

\[
V_g(Q)=\frac12M\omega^2Q^2,
\]

\[
V_e(Q)=E_{eg}+\frac12M\omega^2(Q-\Delta Q)^2.
\]

定义零点涨落

\[
Q_{\rm zpf}=\sqrt{\frac{\hbar}{2M\omega}},
\]

以及无量纲位移

\begin{equation}
\lambda_{\rm FC}
\equiv\frac{\Delta Q}{2Q_{\rm zpf}}
=\sqrt{\frac{M\omega}{2\hbar}}\Delta Q,
\qquad
S_{\rm HR}\equiv\lambda_{\rm FC}^2.
\label{eq:huang-rhys}
\end{equation}

\(S_{\rm HR}\) 即 Huang--Rhys 因子。它可以理解为两个势阱平衡位置间位移相对于零点涨落宽度的平方，也等于经典重组能

\[
E_{\rm reorg}=\frac12M\omega^2(\Delta Q)^2
\]

与一个振动量子 \(\hbar\omega\) 的比值：

\[
S_{\rm HR}=\frac{E_{\rm reorg}}{\hbar\omega}.
\]

Huang--Rhys 因子给出了位移的标度，但 FC 因子的完整结构需用位移算符语言才能看清。以下将激发态振动真空写成基态位移算符作用于基态真空，揭示其相干态本质。

谐振子坐标写成

\[
Q=Q_{\rm zpf}(a+a^\dagger),
\qquad
P=-\ii\sqrt{\frac{M\hbar\omega}{2}}(a-a^\dagger).
\]

平移 \(Q\to Q-\Delta Q\) 的算符为

\[
\exp\left(-\frac{\ii}{\hbar}\Delta Q\,P\right)
=\exp[\lambda_{\rm FC}(a^\dagger-a)].
\]

因此定义位移算符

\begin{equation}
D(\lambda_{\rm FC})
=\exp[\lambda_{\rm FC}(a^\dagger-a)]
\label{eq:displacement-operator-fc}
\end{equation}

（这里 \(\lambda_{\rm FC}\) 取实数）。Baker--Campbell--Hausdorff 公式因为

\[
[\lambda_{\rm FC} a^\dagger,-\lambda_{\rm FC} a]=\lambda_{\rm FC}^2
\]

是 c 数，可精确写成

\begin{equation}
D(\lambda_{\rm FC})
=\ee^{-\lambda_{\rm FC}^2/2}
 \ee^{\lambda_{\rm FC} a^\dagger}
 \ee^{-\lambda_{\rm FC} a}.
\label{eq:displacement-normal-order}
\end{equation}

同时

\[
D^\dagger(\lambda_{\rm FC})aD(\lambda_{\rm FC})=a+\lambda_{\rm FC}.
\]

因此激发态振动真空相对于基态振动基底就是 displaced vacuum，也就是量子光学意义下的相干态。

\begin{proof}
位移算符 \(D(\lambda)=\exp[\lambda(a^\dagger-a)]\)。坐标平移 \(Q\to Q-\Delta Q\) 由生成元 \(P\)（与 \(Q\) 共轭，\([Q,P]=\ii\hbar\)）的指数实现：\(T(\Delta Q)=\exp(-\ii\Delta Q\,P/\hbar)\)。由 BCH 公式 \(\ee^A B\ee^{-A}=B+[A,B]+\cdots\)，取 \(A=\ii\Delta Q\,P/\hbar\)：
\begin{equation*}
T^\dagger Q T=\ee^{\ii\Delta Q\,P/\hbar}\,Q\,\ee^{-\ii\Delta Q\,P/\hbar}=Q+[\tfrac{\ii\Delta Q\,P}{\hbar},Q]+\cdots
\end{equation*}
计算交换子：\([A,Q]=\frac{\ii\Delta Q}{\hbar}[P,Q]=\frac{\ii\Delta Q}{\hbar}(-\ii\hbar)=\Delta Q\)。高阶交换子 \([A,[A,Q]]=[A,\Delta Q]=0\) 截断。因此 \(T^\dagger Q T=Q+\Delta Q\)。代入 \(P=-\ii\sqrt{M\hbar\omega/2}(a-a^\dagger)\)：
\begin{equation*}
T=\exp\!\left(-\frac{\ii\Delta Q}{\hbar}\cdot(-\ii)\sqrt{\frac{M\hbar\omega}{2}}(a-a^\dagger)\right)
=\exp\!\left(-\frac{\Delta Q}{\hbar}\sqrt{\frac{M\hbar\omega}{2}}(a-a^\dagger)\right)
=\exp[\lambda_{\rm FC}(a^\dagger-a)],
\end{equation*}
其中 \(\lambda_{\rm FC}=\Delta Q\sqrt{M\omega/(2\hbar)}=\Delta Q/(2Q_{\rm zpf})\)，即\eqnrefs{eq:huang-rhys}的定义。

令 \(A=\lambda a^\dagger\)，\(B=-\lambda a\)，\([A,B]=\lambda^2\) 为 c 数，BCH 公式 \(\ee^{A+B}=\ee^A\ee^B\ee^{-[A,B]/2}\) 给出
\begin{equation*}
D(\lambda)=\ee^{\lambda(a^\dagger-a)}=\ee^{\lambda a^\dagger}\ee^{-\lambda a}\ee^{-\lambda^2/2},
\end{equation*}
即\eqnrefs{eq:displacement-normal-order}。

由 \([a,a^\dagger]=1\)，\([a,\lambda a^\dagger]=\lambda\)，\([a,-\lambda a]=0\)。BCH 公式：
\begin{equation*}
D^\dagger(\lambda)aD(\lambda)=\ee^{-\lambda(a^\dagger-a)}\,a\,\ee^{\lambda(a^\dagger-a)}
=a+[-\lambda(a^\dagger-a),a]+\frac{1}{2}[-\lambda(a^\dagger-a),[-\lambda(a^\dagger-a),a]]+\cdots
\end{equation*}
指数左侧的生成元为 $-\lambda(a^\dagger-a)$，其与 $a$ 的交换子为 $+\lambda$；这个数再取交换子为零，所以展开在一阶截断。因此
\begin{equation*}
D^\dagger a D=a+\lambda.
\end{equation*}
即 \(a(D|0\rangle)=\lambda(D|0\rangle)\)。位移算符 \(D(\lambda_{\rm FC})\) 把基态振动真空平移到激发态振动真空，\(\lambda_{\rm FC}=\Delta Q/(2Q_{\rm zpf})\) 即两势阱平衡位置之差以零点宽度为单位，FC 因子由 Huang--Rhys 因子 \(S_{\rm HR}=\lambda_{\rm FC}^2\) 决定。综合以上，位移算符的关键代数性质为
\begin{equation*}
D^\dagger(\lambda)\,a\,D(\lambda)=a+\lambda,\qquad D^\dagger(\lambda)\,a^\dagger\,D(\lambda)=a^\dagger+\lambda^*.
\end{equation*}

\end{proof}

有了真空-真空重叠，一般 $m\to n$ 跃迁的 FC 因子只需在位移算符两侧插入升降算符。以下用生成函数方法一次性获得全部 FC 因子的封闭表达式。

设初态为基态势阱的振动真空。由于

\[
\ee^{-\lambda_{\rm FC} a}|0\rangle=|0\rangle,
\]

有

\[
D(\lambda_{\rm FC})|0\rangle
=\ee^{-\lambda_{\rm FC}^2/2}
\sum_{n=0}^{\infty}
\frac{\lambda_{\rm FC}^n}{n!}(a^\dagger)^n|0\rangle.
\]

利用

\[
(a^\dagger)^n|0\rangle=\sqrt{n!}|n\rangle,
\]

得到

\begin{equation}
\langle n|D(\lambda_{\rm FC})|0\rangle
=\ee^{-\lambda_{\rm FC}^2/2}\frac{\lambda_{\rm FC}^n}{\sqrt{n!}}.
\label{eq:fc-amplitude-ground}
\end{equation}

因此

\begin{equation}
|\langle n_e|0_g\rangle|^2
=\ee^{-S_{\rm HR}}
\frac{S_{\rm HR}^n}{n!}.
\label{eq:fc-poisson}
\end{equation}

不同 Huang--Rhys 因子下的分布见\figref{fig:franck-condon}。

\begin{figure}[tbp]
\centering
\includegraphics[width=.83\textwidth]{franck_condon.pdf}
\caption{等频率位移谐振子模型中的 Franck--Condon 因子。Huang--Rhys 因子 \(S_{\rm HR}\) 同时决定 Poisson 分布的均值和方差。}
\label{fig:franck-condon}
\end{figure}

\eqnrefs{eq:fc-poisson} 的均值和方差都是 \(S_{\rm HR}\)。因此 \(S_{\rm HR}\ll1\) 时零声子线占主导；\(S_{\rm HR}\gg1\) 时强度分布中心移向 \(n\sim S_{\rm HR}\)，出现明显振动 progression。

\begin{proposition}[任意振动态之间的位移重叠]
在等频率位移谐振子模型中，实位移参数 $\lambda_{\rm FC}$ 给出
\[
\langle m|D(\lambda_{\rm FC})|n\rangle
=\ee^{-\lambda_{\rm FC}^2/2}
\begin{cases}
\sqrt{n!/m!}\,\lambda_{\rm FC}^{m-n}
L_n^{(m-n)}(\lambda_{\rm FC}^2),&m\ge n,\\
\sqrt{m!/n!}\,(-\lambda_{\rm FC})^{n-m}
L_m^{(n-m)}(\lambda_{\rm FC}^2),&n\ge m.
\end{cases}
\]
两种写法在 $m=n$ 时一致。Franck--Condon 因子是该振幅的模平方；
位移的符号虽不改变单个跃迁概率，却影响振幅的相对相位。
\end{proposition}
\begin{proof}[证明（有限求和与 Laguerre 多项式）]
先取 $m\ge n$。对任意数态，使用 \eqnrefs{eq:displacement-normal-order}：
\begin{equation*}
\langle m|D(\lambda_{\rm FC})|n\rangle
=\ee^{-\lambda_{\rm FC}^2/2}
\langle m|\ee^{\lambda_{\rm FC} a^\dagger}\ee^{-\lambda_{\rm FC} a}|n\rangle.
\end{equation*}
先展开右侧的 \(\ee^{-\lambda_{\rm FC} a}\)，利用 \(a^j|n\rangle=\sqrt{n!/(n-j)!}\,|n-j\rangle\) 且 \(a^j|n\rangle=0\) 对 \(j>n\)，级数自动截断：
\begin{equation*}
\ee^{-\lambda_{\rm FC} a}|n\rangle
=\sum_{j=0}^{\infty}\frac{(-\lambda_{\rm FC})^j}{j!}a^j|n\rangle
=\sum_{j=0}^{n}\frac{(-\lambda_{\rm FC})^j}{j!}
\sqrt{\frac{n!}{(n-j)!}}\,|n-j\rangle.
\end{equation*}
再展开左侧的 \(\ee^{\lambda_{\rm FC} a^\dagger}\) 并利用数态正交性 \(\langle m|a^{\dagger i}|n-j\rangle=\sqrt{m!/(m-i)!}\,\delta_{m-i,n-j}\)，得到
\begin{align*}
\langle m|D(\lambda_{\rm FC})|n\rangle
&=\ee^{-\lambda_{\rm FC}^2/2}
\sum_{j=0}^{n}\frac{(-\lambda_{\rm FC})^j}{j!}
\sqrt{\frac{n!}{(n-j)!}}
\sum_{i=0}^{m}\frac{\lambda_{\rm FC}^i}{i!}
\sqrt{\frac{m!}{(m-i)!}}\,\delta_{m-i,n-j}\\
&=\ee^{-\lambda_{\rm FC}^2/2}\sqrt{n!\,m!}
\sum_{j=0}^{n}\frac{(-1)^j\lambda_{\rm FC}^{m-n+2j}}{(m-n+j)!\,j!\,(n-j)!},
\end{align*}
其中 Kronecker 记号固定 \(i=m-n+j\)，故求和指标仍为 \(j\)，且要求 \(m\ge n\)。把 \(\lambda_{\rm FC}^{m-n+2j}=\lambda_{\rm FC}^{m-n}\lambda_{\rm FC}^{2j}\) 拆开并提出求和号外，引入关联 Laguerre 多项式
\begin{equation*}
L_n^{(\alpha)}(x)
=\sum_{j=0}^{n}\frac{(-x)^j}{j!}\,
\frac{\Gamma(n+\alpha+1)}{\Gamma(\alpha+j+1)\,(n-j)!},
\end{equation*}
取 \(\alpha=m-n\) 时 \(\Gamma(m+1)=m!\)、\(\Gamma(m-n+j+1)=(m-n+j)!\)，于是求和正好满足
\begin{equation*}
\sum_{j=0}^{n}\frac{(-1)^j\lambda_{\rm FC}^{m-n+2j}}{(m-n+j)!\,j!\,(n-j)!}
=\frac{\lambda_{\rm FC}^{m-n}}{m!}
L_n^{(m-n)}(\lambda_{\rm FC}^2),
\end{equation*}
代回前式并约去 \(m!\)，
\begin{align*}
\langle m|D(\lambda_{\rm FC})|n\rangle
&=\ee^{-\lambda_{\rm FC}^2/2}\sqrt{n!\,m!}\,
\frac{\lambda_{\rm FC}^{m-n}}{m!}
L_n^{(m-n)}(\lambda_{\rm FC}^2)\\
&=\ee^{-\lambda_{\rm FC}^2/2}
\sqrt{\frac{n!}{m!}}\,\lambda_{\rm FC}^{m-n}
L_n^{(m-n)}(\lambda_{\rm FC}^2).
\end{align*}
因此对于 \(m\ge n\)，
\begin{equation*}
\langle m|D(\lambda_{\rm FC})|n\rangle
=\ee^{-\lambda_{\rm FC}^2/2}
\sqrt{\frac{n!}{m!}}
\lambda_{\rm FC}^{m-n}
L_n^{(m-n)}(\lambda_{\rm FC}^2),
\end{equation*}
对于 $n\ge m$，不能简单交换 $m,n$ 而不改变位移符号。
由 $D(\lambda)^\dagger=D(-\lambda)$，
\[
\langle m|D(\lambda)|n\rangle
=\langle n|D(-\lambda)|m\rangle^*,
\]
代入已证明的第一种情形，得到含 $(-\lambda)^{n-m}$ 的第二种公式。
$n=0$ 时 $L_0^{(m)}=1$，恢复\eqnrefs{eq:fc-amplitude-ground}。
\end{proof}

\begin{example}[激发振动态的节点与强度和]
若初态为 $\ket1$，令 $S=\lambda_{\rm FC}^2$，则
\[
|\langle0|D|1\rangle|^2=\ee^{-S}S,\qquad
|\langle1|D|1\rangle|^2=\ee^{-S}(1-S)^2.
\]
第二条跃迁在 $S=1$ 处消失，是 Laguerre 多项式的零点；
激发态的强度分布因此不再是基态的 Poisson 分布。
完整振动态基与 $D$ 的酉性给出强度和
\[
\sum_{m=0}^\infty |\langle m|D|n\rangle|^2
=\langle n|D^\dagger D|n\rangle=1.
\]
若截断高振动态，这个和小于一，差额表示截断掉的跃迁权重。
\end{example}


等频率位移只涉及位移算符 $D(\alpha)$。但实际分子的两个电子态势能面通常频率不同、正规模方向也不一致，分别引入压缩和 Duschinsky 旋转。本节处理这一一般情形。

\subsection{Duschinsky--Doktorov 变换}
如果两电子态的振动频率不同，

\[
\omega_g\neq\omega_e,
\]

即使平衡位置相同，两组升降算符也不同。用同一组物理 \(Q,P\) 写

\[
a_g
=\sqrt{\frac{M\omega_g}{2\hbar}}Q
+\frac{\ii}{\sqrt{2M\hbar\omega_g}}P,
\]

\[
a_e
=\sqrt{\frac{M\omega_e}{2\hbar}}Q
+\frac{\ii}{\sqrt{2M\hbar\omega_e}}P.
\]

消去 \(Q,P\)，得到 Bogoliubov 变换

\[
a_e
=u a_g+v a_g^\dagger,
\]

\[
u=\frac12\left(
\sqrt{\frac{\omega_e}{\omega_g}}
+\sqrt{\frac{\omega_g}{\omega_e}}
\right),
\qquad
v=\frac12\left(
\sqrt{\frac{\omega_e}{\omega_g}}
-\sqrt{\frac{\omega_g}{\omega_e}}
\right),
\]

满足 \(u^2-v^2=1\)。令

\[
r=\frac12\ln\frac{\omega_e}{\omega_g},
\]

则 \(u=\cosh r,v=\sinh r\)。所以“频率改变对应压缩”不是类比，而是两套谐振子正则坐标之间的精确 Bogoliubov 关系。

频率改变只引入对角压缩。多原子分子中两个电子态的正规坐标方向一般不同，还需一个正交旋转——Duschinsky 变换——来描述模式混合。

多原子分子的两个电子态一般具有不同 Hessian，因此正规模方向会旋转。采用质量加权正规坐标，最常见的线性关系写成

\begin{equation}
\mathbf Q_e
=\mathbf J\mathbf Q_g+\mathbf K,
\label{eq:duschinsky}
\end{equation}

其中 \(\mathbf J\) 是 Duschinsky 转动，\(\mathbf K\) 是平衡构型位移。若两组质量加权坐标都正交归一，则 \(\mathbf J\) 为正交矩阵，并且为保持正则结构

\[
\mathbf P_e=\mathbf J\mathbf P_g.
\]

令 \(\boldsymbol\Omega_g\)、\(\boldsymbol\Omega_e\) 分别是由各振动频率组成的对角矩阵，定义多模升降算符

\[
\mathbf a_s
=\frac1{\sqrt{2\hbar}}
\left(
\boldsymbol\Omega_s^{1/2}\mathbf Q_s
+\ii\boldsymbol\Omega_s^{-1/2}\mathbf P_s
\right),
\qquad s=g,e.
\]

把 \eqnrefs{eq:duschinsky} 代入后得到一般高斯仿射变换

\begin{equation}
\mathbf a_e
=\mathbf X\mathbf a_g
+\mathbf Y\mathbf a_g^\dagger
+\mathbf z,
\label{eq:doktorov-bogoliubov}
\end{equation}

其中

\[
\mathbf X
=\frac12
\left(
\boldsymbol\Omega_e^{1/2}\mathbf J\boldsymbol\Omega_g^{-1/2}
+
\boldsymbol\Omega_e^{-1/2}\mathbf J\boldsymbol\Omega_g^{1/2}
\right),
\]

\[
\mathbf Y
=\frac12
\left(
\boldsymbol\Omega_e^{1/2}\mathbf J\boldsymbol\Omega_g^{-1/2}
-
\boldsymbol\Omega_e^{-1/2}\mathbf J\boldsymbol\Omega_g^{1/2}
\right),
\]

\[
\mathbf z
=\frac1{\sqrt{2\hbar}}
\boldsymbol\Omega_e^{1/2}\mathbf K.
\]

多模 Franck--Condon 变换因而保持玻色对易关系：\(\mathbf J\) 负责模式混合，\(\mathbf X,\mathbf Y\) 中的频率比负责压缩，\(\mathbf z\) 负责位移。它们合起来称为 Doktorov 变换。

实际大分子的 Franck--Condon 因子因而是多维高斯波函数之间的重叠问题；模式旋转和压缩会让不同简正模之间的量子数不再逐模独立守恒。

以上均在零温（初态为振动基态）下讨论。有限温度下初态为热态，FC 因子需对热平均求和，用时间关联函数可以给出紧凑的表达。

\subsection{有限温度谱与时间关联函数}

在有限温度下，初始振动态不再是单个 $|0_g\rangle$，而是热态
\[
\rho_g^{\rm th}=\frac{\ee^{-\beta H_g}}{Z_g},
\qquad \beta=(k_BT)^{-1}.
\]
Condon 近似下，吸收谱可以避免逐个枚举所有初末振动态，而写成偶极关联函数的 Fourier 变换
\begin{equation}
I(\omega)\propto
\int_{-\infty}^{\infty}\dd t\,
\ee^{\ii\omega t}
C_{eg}(t),
\qquad
C_{eg}(t)=\operatorname{Tr}\!\left[
ho_g^{\rm th}
\ee^{\ii H_gt/\hbar}\ee^{-\ii H_et/\hbar}\right].
\label{eq:franck-condon-thermal-correlation}
\end{equation}
对一般二次 $H_g,H_e$，\eqnrefs{eq:doktorov-bogoliubov}保证这个迹仍是高斯迹，因此可以由协方差矩阵和位移向量闭合计算。零温、等频率、单模位移模型只是在该一般生成函数中取 $T\to0$、$\mathbf J=I$、$\mathbf Y=0$ 后得到的 Poisson 特例。

\begin{proof}
离散振动态的吸收强度由 Fermi 黄金法则给出
\begin{equation*}
I(\omega)\propto
\sum_{m,n} p_n\,
|\langle e,m|\hat\mu|g,n\rangle|^2\,
\delta\!\left(\omega-\frac{E_m^e-E_n^g}{\hbar}\right),
\end{equation*}
其中 $p_n=\ee^{-\beta E_n^g}/Z_g$ 为初态 Boltzmann 权重，$|g,n\rangle$、$|e,m\rangle$ 分别为电子基态与激发态势能面上的振动态。记 $\omega_{mn}=(E_m^e-E_n^g)/\hbar$，将矩阵元模平方展开，
\begin{equation*}
|\langle e,m|\hat\mu|g,n\rangle|^2
= \langle e,m|\hat\mu|g,n\rangle\,
 \langle g,n|\hat\mu^\dagger|e,m\rangle.
\end{equation*}
利用 $\delta$ 函数的 Fourier 表示
\begin{equation*}
\delta(\omega-\omega_{mn})
= \frac1{2\pi}\int_{-\infty}^{\infty}\dd t\,
 \ee^{\ii(\omega-\omega_{mn})t},
\end{equation*}
代回强度公式并交换求和与积分，
\begin{equation*}
I(\omega)\propto
\frac1{2\pi}\int_{-\infty}^{\infty}\dd t\,
\ee^{\ii\omega t}
\sum_{m,n} p_n\,
\langle e,m|\hat\mu|g,n\rangle
\langle g,n|\hat\mu^\dagger|e,m\rangle
\ee^{-\ii\omega_{mn}t}.
\end{equation*}
由 $H_g|g,n\rangle=E_n^g|g,n\rangle$、$H_e|e,m\rangle=E_m^e|e,m\rangle$，于是
\begin{equation*}
\ee^{-\ii\omega_{mn}t}
= \langle e,m|\ee^{-\ii H_e t/\hbar}|e,m\rangle\,
 \langle g,n|\ee^{+\ii H_g t/\hbar}|g,n\rangle.
\end{equation*}
在 Condon 近似下电子偶极与核坐标无关，$\hat\mu=\boldsymbol\mu_{eg}\otimes\hat{\mathbf 1}_{\rm vib}$，可提出电子因子。对 $m$ 与 $n$ 分别插入完备关系 $\sum_m|e,m\rangle\langle e,m|=\mathbf 1$、$\sum_n|g,n\rangle\langle g,n|=\mathbf 1$，双重求和坍缩为
\begin{equation*}
\sum_n p_n\,
\langle g,n|\ee^{+\ii H_g t/\hbar}\ee^{-\ii H_e t/\hbar}|g,n\rangle
= \operatorname{Tr}\!\left[
 \rho_g^{\rm th}\,
 \ee^{+\ii H_g t/\hbar}\ee^{-\ii H_e t/\hbar}\right],
\end{equation*}
其中 $\rho_g^{\rm th}=\sum_n p_n|g,n\rangle\langle g,n|=\ee^{-\beta H_g}/Z_g$。因此
\begin{equation*}
I(\omega)\propto
\int_{-\infty}^{\infty}\dd t\,
\ee^{\ii\omega t}\,
C_{eg}(t),
\qquad
C_{eg}(t)=\operatorname{Tr}\!\left[
\rho_g^{\rm th}\,
\ee^{+\ii H_g t/\hbar}\ee^{-\ii H_e t/\hbar}\right],
\end{equation*}
即\eqnrefs{eq:franck-condon-thermal-correlation}。大分子振动谱更自然的母对象是高斯关联函数，而非数量指数增长的逐态 Franck--Condon 表。
\end{proof}

\subsection{Franck--Condon 因子的实验数值与典型体系}\label{sec:fc-experiments}

Huang--Rhys 因子、Duschinsky 旋转和有限温度关联函数分别控制位移、模式混合与热布居。代入具体分子的振动频率和平衡构型差以后，它们直接给出谱带包络和各振动支线的相对强度。

单模位移谐振子的 Huang--Rhys 因子 \(S_{\rm HR}=(\Delta Q)^2/(2\ell_v^2)\) 决定振动电子谱线强度分布：\(I_n=e^{-S_{\rm HR}}S_{\rm HR}^{\,n}/n!\)（Poisson 分布）。\(S_{\rm HR}\) 的物理意义是基态到末态振动波包中心的位移（以零点宽度为单位）。

\begin{center}
\begin{tabular}{lcccc}
\hline
体系 & 跃迁 & \(S\) & Stokes 位移 (cm\(^{-1}\)) & \(\lambda_{\rm FC}^2\) \\
\hline
H\(_2\) (\(B\,{}^1\Sigma\leftarrow X\,{}^1\Sigma\)) & Lyman & 0.09 & 1700 & 0.18 \\
I\(_2\) (\(B\,{}^3\Pi\leftarrow X\,{}^1\Sigma\)) & 可见 & 1.5 & 240 & 3.0 \\
F-center (NaCl) & \(1s\to2p\) & 3.0 & 1800 & 6.0 \\
聚芴 (有机半导体) & \(S_1\to S_0\) & 2.0 & 2400 & 4.0 \\
CdSe 量子点 & 带边 & 0.5 & 200 & 1.0 \\
\hline
\end{tabular}
\end{center}

\begin{exbox}[\(I_2\) 的振动电子进展]
I\(_2\) 的 \(B\leftarrow X\) 跃迁 \(S\approx1.5\)，振动电子峰强分布 \(I_n=e^{-1.5}1.5^n/n!\)：\(I_0=0.22\)，\(I_1=0.34\)，\(I_2=0.25\)，\(I_3=0.13\)，\(I_4=0.05\)。最强峰在 \(n\approx S=1.5\) 附近，进展宽度 \(\Delta n\approx\sqrt{S}\approx1.2\)。对比 H\(_2\) 的 \(S=0.09\)：\(I_0=0.91\)，\(I_1=0.08\)——几乎只有 0-0 带，因为 H\(_2\) 两个电子态的平衡键长几乎相同。
\end{exbox}

Huang--Rhys 因子决定振动电子包络的形状，但光谱中哪些跃迁实际出现还受选择定则约束。大 $S$ 时反射原理将 FC 因子与连续谱阈联系起来。

\begin{exbox}[反射原理：\(I_2\) 吸收谱的连续阈]
当 \(S\gg1\) 且末态势能面显著位移时，FC 因子在 \(v'\approx S\) 处达到峰值，对应经典转折点。\(I_2\) 的 \(B\) 态势能在 \(R\approx3.0\,\text{\AA}\) 处有势阱，而基态 \(R_e=2.67\,\text{\AA}\)。当激发光能量超过 \(B\) 态离解限时，FC 因子给出连续吸收谱——反射原理：核波函数在转折点反射，吸收谱形 \(\sigma(\nu)\propto|\chi_g(R_t)|^2/|V_e'(R_t)-V_g'(R_t)|\)，其中 \(R_t\) 满足 \(V_e(R_t)-V_g(R_t)=h\nu\)。
\end{exbox}

以上实验体系均为单模或频率一致情形。一般多原子分子中 Duschinsky 旋转使模式混合，FC 因子成为多维重叠积分。

当两个电子态的正规模方向不同时，Duschinsky 变换 \(Q'=UQ+\Delta Q\) 引入模式混合。同一对称构型下的模式混合还受点群表示约束：属于不同不可约表示的模式不会仅因改变频率而相互转化。实际构型改变或对称性降低时，才须在共同坐标空间中重新计算允许的混合矩阵。下面用明确给定的谐振子模型计算重叠，避免将模型参数当作未注明测量条件的分子常数。

\subsection{多模重叠与高斯积分的完整计算}

电子跃迁涉及两套谐振子，但核坐标仍是同一个空间中的变量。以下统一使用质量加权坐标 $\mathbf Q\in\mathbb R^N$，动能为 $\mathbf P^{\mathsf T}\mathbf P/2$；若使用有长度量纲的单模坐标，先作 $Q=\sqrt M\,q$。设两套频率矩阵 $\boldsymbol\Omega_g,\boldsymbol\Omega_e$ 均为严格正的实对角矩阵，并取
\begin{equation}
 \mathbf Q_e=U\mathbf Q+\mathbf d,\qquad
 \mathbf P_e=U\mathbf P,\qquad U^{\mathsf T}U=I.
 \label{eq:duschinsky-general}
\end{equation}
这与\eqnrefs{eq:duschinsky}中的 $U=\mathbf J$、$\mathbf d=\mathbf K$ 为同一约定。两个核哈密顿量写成
\[
 H_g=\tfrac12\mathbf P^{\mathsf T}\mathbf P
       +\tfrac12\mathbf Q^{\mathsf T}\boldsymbol\Omega_g^2\mathbf Q,
 \qquad
 H_e=E_{eg}+\tfrac12\mathbf P_e^{\mathsf T}\mathbf P_e
       +\tfrac12\mathbf Q_e^{\mathsf T}\boldsymbol\Omega_e^2\mathbf Q_e.
\]
这里 $E_{eg}$ 是两个势阱最低点的电子能量差。末态平衡位置在初态坐标中为 $-U^{\mathsf T}\mathbf d$。全部模频率为正，使核热态有有限配分函数；平动、转动零模须在建立这套振动模型之前分离。

\begin{proposition}[模式旋转的正则条件]
上述坐标变换保持 $[Q_i,P_j]=\ii\hbar\delta_{ij}$。相应升降算符的系数满足
\begin{equation}
 XX^{\mathsf T}-YY^{\mathsf T}=I,\qquad
 XY^{\mathsf T}-YX^{\mathsf T}=0.
 \label{eq:duschinsky-symplectic}
\end{equation}
这里 $X,Y$ 是\eqnrefs{eq:doktorov-bogoliubov}中给出的实矩阵；频率比本身无需保持一个频率加权内积。
\end{proposition}
\begin{proof}
坐标的交换子为 $[Q_{e,i},P_{e,j}]=\ii\hbar(UU^{\mathsf T})_{ij}=\ii\hbar\delta_{ij}$。为了核对升降算符中的排序，记
\[
 R=\boldsymbol\Omega_e^{1/2}U\boldsymbol\Omega_g^{-1/2},
 \qquad S=\boldsymbol\Omega_e^{-1/2}U\boldsymbol\Omega_g^{1/2},
 \qquad X=(R+S)/2,\quad Y=(R-S)/2.
\]
正交性给出 $RS^{\mathsf T}=SR^{\mathsf T}=I$，因而
\[
 XX^{\mathsf T}-YY^{\mathsf T}
 =\tfrac12(RS^{\mathsf T}+SR^{\mathsf T})=I,
 \qquad
 XY^{\mathsf T}-YX^{\mathsf T}
 =\tfrac12(SR^{\mathsf T}-RS^{\mathsf T})=0.
\]
把 $a_e=Xa_g+Ya_g^\dagger+z$ 代入交换子，就分别得到 $[a_{e,i},a_{e,j}^\dagger]=\delta_{ij}$ 与 $[a_{e,i},a_{e,j}]=0$；位移 $z$ 是数，不改变交换子。
\end{proof}

先计算振幅，而不是直接平方。实振动本征函数的重叠为
\begin{equation}
 F_{\mathbf m,\mathbf n}
 =\int_{\mathbb R^N}\dd^NQ\,
 \chi_{e,\mathbf m}(U\mathbf Q+\mathbf d)^*
 \chi_{g,\mathbf n}(\mathbf Q).
 \label{eq:multimode-fc-integral}
\end{equation}
保留振幅的符号，才能在有多条光学路径时继续计算干涉。令 $\mathbf z,\mathbf z'$ 为独立的形式变量，定义
\begin{equation}
 G(\mathbf z,\mathbf z')
 =\sum_{\mathbf n,\mathbf m}F_{\mathbf m,\mathbf n}
   \prod_i\frac{z_i^{n_i}}{\sqrt{n_i!}}
   \prod_j\frac{{z'_j}^{m_j}}{\sqrt{m_j!}}.
 \label{eq:multimode-fc-amplitude-generator}
\end{equation}

\begin{proposition}[带归一化的振幅生成函数]
记
\[
 L=\boldsymbol\Omega_g+U^{\mathsf T}\boldsymbol\Omega_eU,
 \qquad h=U^{\mathsf T}\boldsymbol\Omega_e\mathbf d.
\]
则
\begin{equation}
 G=F_{\mathbf0,\mathbf0}
 \exp\!\left[
 \tfrac12\mathbf z^{\mathsf T}A\mathbf z
 +\tfrac12{\mathbf z'}^{\mathsf T}B\mathbf z'
 +\mathbf z^{\mathsf T}C\mathbf z'
 +l^{\mathsf T}\mathbf z+{l'}^{\mathsf T}\mathbf z'
 \right],
 \label{eq:doktorov-gaussian-generator}
\end{equation}
其中真空重叠和全部系数为
\begin{align*}
 F_{\mathbf0,\mathbf0}
 &=\frac{2^{N/2}\det(\boldsymbol\Omega_g\boldsymbol\Omega_e)^{1/4}}
         {\sqrt{\det L}}
   \exp\!\left[-\frac{\mathbf d^{\mathsf T}\boldsymbol\Omega_e\mathbf d
                      -h^{\mathsf T}L^{-1}h}{2\hbar}\right],\\
 A&=2\boldsymbol\Omega_g^{1/2}L^{-1}\boldsymbol\Omega_g^{1/2}-I,\\
 B&=2\boldsymbol\Omega_e^{1/2}UL^{-1}U^{\mathsf T}\boldsymbol\Omega_e^{1/2}-I,\\
 C&=2\boldsymbol\Omega_g^{1/2}L^{-1}U^{\mathsf T}\boldsymbol\Omega_e^{1/2},\\
 l&=-\sqrt{2/\hbar}\,\boldsymbol\Omega_g^{1/2}L^{-1}h,\\
 l'&=\sqrt{2/\hbar}\,\boldsymbol\Omega_e^{1/2}(\mathbf d-UL^{-1}h).
\end{align*}
正定的 $L$ 使平方根在实正支上唯一确定。
\end{proposition}
\begin{proof}
对质量加权谐振子，Hermite 多项式的生成函数给出
\[
 \sum_{\mathbf n}\chi_{g,\mathbf n}(\mathbf Q)
 \prod_i\frac{z_i^{n_i}}{\sqrt{n_i!}}
 =\chi_{g,\mathbf0}(\mathbf Q)
 \exp\!\left[\sqrt{2/\hbar}\,
 \mathbf z^{\mathsf T}\boldsymbol\Omega_g^{1/2}\mathbf Q
 -\tfrac12\mathbf z^{\mathsf T}\mathbf z\right],
\]
其中 $\chi_{g,\mathbf0}=\det(\boldsymbol\Omega_g/(\pi\hbar))^{1/4}
\exp[-\mathbf Q^{\mathsf T}\boldsymbol\Omega_g\mathbf Q/(2\hbar)]$。
末态使用同一公式，将 $\mathbf Q$ 换成 $U\mathbf Q+\mathbf d$、$\mathbf z$ 换成 $\mathbf z'$。

把两个生成函数乘入重叠积分。令
$v=\boldsymbol\Omega_g^{1/2}\mathbf z+U^{\mathsf T}\boldsymbol\Omega_e^{1/2}\mathbf z'$，关于 $\mathbf Q$ 的指数为
\[
 -\frac1{2\hbar}\mathbf Q^{\mathsf T}L\mathbf Q
 +\frac1\hbar(\sqrt{2\hbar}\,v-h)^{\mathsf T}\mathbf Q.
\]
完成平方后，积分贡献 $(2\pi\hbar)^{N/2}/\sqrt{\det L}$，以及
$(\sqrt{2\hbar}\,v-h)^{\mathsf T}L^{-1}(\sqrt{2\hbar}\,v-h)/(2\hbar)$。
常数项给出 $F_{\mathbf0,\mathbf0}$，二次项与一次项分别给出 $A,B,C,l,l'$。因此归一化因子与位移指数都由同一个高斯积分确定。
\end{proof}

对\eqnrefs{eq:doktorov-gaussian-generator}求 $z_i$ 导数，并逐项比较幂级数系数，得到可直接使用的递推：
\begin{equation}
\begin{aligned}
 \sqrt{n_i+1}\,F_{\mathbf m,\mathbf n+\mathbf e_i}
 ={}&\sum_j A_{ij}\sqrt{n_j}\,F_{\mathbf m,\mathbf n-\mathbf e_j}\\
 &+\sum_j C_{ij}\sqrt{m_j}\,F_{\mathbf m-\mathbf e_j,\mathbf n}
 +l_iF_{\mathbf m,\mathbf n}.
\end{aligned}
\label{eq:doktorov-fc-recurrence}
\end{equation}
负量子数项取零。对 $z'_i$ 求导可得交换初末态的另一递推；从已知的 $F_{\mathbf0,\mathbf0}$ 出发，先生成 $\mathbf n=0$ 的末态振幅，再使用上式增加初态量子数。

\begin{workedexample}[位移与频率改变怎样共同改变真空重叠]
取单模 $U=1$。求任意正频率 $\omega_g,\omega_e$ 与质量加权位移 $d$ 下的真空重叠，并计算等频率时的 $0\to0$、$0\to1$ 振幅。

\textbf{解。}代入 $L=\omega_g+\omega_e$、$h=\omega_ed$，得到
\[
 F_{0,0}
 =\left(\frac{2\sqrt{\omega_g\omega_e}}{\omega_g+\omega_e}\right)^{1/2}
 \exp\!\left[-\frac{\omega_g\omega_ed^2}
 {2\hbar(\omega_g+\omega_e)}\right].
\]
无位移而 $\omega_e=4\omega_g$ 时，$|F_{0,0}|^2=4/5$；频率改变本身已把其余权重分配到激发态。等频率 $\omega_g=\omega_e=\omega$ 时，$A=B=0$、$C=1$、$l=-\lambda$、$l'=\lambda$，其中 $\lambda=\sqrt{\omega/(2\hbar)}d$，于是
\[
 G=\ee^{-\lambda^2/2}\exp(zz'-\lambda z+\lambda z'),
 \qquad F_{0,0}=\ee^{-\lambda^2/2},\quad
 F_{1,0}=\lambda\ee^{-\lambda^2/2}.
\]
这同时核对了生成函数的常数、位移符号和\eqnrefs{eq:fc-poisson}的概率。更一般的单模重叠即前面已证明的 Laguerre 公式，可保留同一标签写为
\begin{equation}
 F_{m,n}=\langle m|D(\lambda)|n\rangle,
 \qquad \lambda=\sqrt{\omega/(2\hbar)}d.
 \label{eq:huang-rhys-laguerre-overlap}
\end{equation}
交换 $m,n$ 时还须按 $D^\dagger(\lambda)=D(-\lambda)$ 改变位移符号。
\end{workedexample}

\subsection{有限温度关联的可计算高斯形式}

振幅生成函数解决逐态重叠。若只求整个有限温度谱，\eqnrefs{eq:franck-condon-thermal-correlation}中的迹可以一次算出。这里沿用吸收谱的正号 Fourier 约定，令
\begin{equation}
 C(t)=Z_g^{-1}\Tr\!\left[\ee^{-(\beta-\ii t/\hbar)H_g}
                              \ee^{-\ii H_et/\hbar}\right],
 \qquad Z_g=\prod_j[2\sinh(\beta\hbar\omega_{g,j}/2)]^{-1}.
 \label{eq:thermal-correlation}
\end{equation}
因此 $C(0)=1$，谱的全部振动强度之和也为一。有限温度会同时保留升振动量子数与降振动量子数的跃迁。

\begin{proposition}[一般双谐振势能面的热关联]
取 $0<\eta<\beta$，令 $s_g=\beta-\eta-\ii t/\hbar$、$s_e=\eta+\ii t/\hbar$。对 $a=g,e$ 定义
\begin{equation}
 A_a=\boldsymbol\Omega_a\coth(\hbar s_a\boldsymbol\Omega_a),
 \qquad B_a=\boldsymbol\Omega_a\operatorname{csch}(\hbar s_a\boldsymbol\Omega_a).
 \label{eq:coherent-state-propagator}
\end{equation}
这些是频率对角矩阵的函数。再记
\[
 D_t=A_g+U^{\mathsf T}A_eU,\qquad
 F_t=B_g+U^{\mathsf T}B_eU,\qquad
 M_t=\begin{pmatrix}D_t&-F_t\\-F_t&D_t\end{pmatrix},
\]
\[
 h_t=\begin{pmatrix}U^{\mathsf T}(A_e-B_e)\mathbf d\\
                     U^{\mathsf T}(A_e-B_e)\mathbf d\end{pmatrix},
 \qquad c_t=2\mathbf d^{\mathsf T}(A_e-B_e)\mathbf d.
\]
则\eqnrefs{eq:thermal-correlation}等于
\begin{equation}
 C(t)=\lim_{\eta\downarrow0}
 \frac{\ee^{-s_eE_{eg}}}{Z_g}
 \left[\frac{\det B_g\det B_e}{\det M_t}\right]^{1/2}
 \exp\!\left[\frac{h_t^{\mathsf T}M_t^{-1}h_t-c_t}{2\hbar}\right].
 \label{eq:multimode-fc-generating-function}
\end{equation}
对每个固定的 $0<\eta<\beta$，平方根从 $t=0$ 的正热核沿实 $t$ 连续延拓，再取 $\eta\downarrow0$。不能在各个时刻分别选取复数的主平方根，否则经过相位绕转后会改变关联函数的符号。单独 Mehler 因子在实时间极限中的焦散，由整个表达的连续极限处理。
\end{proposition}
\begin{proof}
数学输入是已证明的 Mehler 核\eqnrefs{eq:phys-ghs-mehler-kernel}。对 $\Re s>0$，它在质量加权坐标中的多模乘积为
\[
 K_a(q,q';s)=\frac{\det(B_a)^{1/2}}{(2\pi\hbar)^{N/2}}
 \exp\!\left[-\frac{q^{\mathsf T}A_aq+{q'}^{\mathsf T}A_aq'
                         -2q^{\mathsf T}B_aq'}{2\hbar}\right].
\]
取迹时插入两次坐标完备关系，得到
\[
 C_\eta(t)=\frac{\ee^{-s_eE_{eg}}}{Z_g}
 \int\dd^Nq\dd^Nq'\,
 K_g(q,q';s_g)K_e(Uq'+\mathbf d,Uq+\mathbf d;s_e).
\]
因为 $U$ 正交，坐标变换的 Jacobian 绝对值为一。将两个指数相加，令 $x=(q,q')^{\mathsf T}$，得到
\[
 -\frac1{2\hbar}(x^{\mathsf T}M_tx+2h_t^{\mathsf T}x+c_t).
\]
完成平方后，$2N$ 维积分贡献 $(2\pi\hbar)^N\det(M_t)^{-1/2}$；它与两个核的前因子约去，给出所述行列式与位移指数。$\eta>0$ 使两个核都是有收敛因子的热核，随后取 $\eta\downarrow0$ 才恢复实时间演化。

正频率使两个哈密顿量下有界。对 $0<\eta<\beta/2$，在初态数态基中，迹的第 $\mathbf n$ 项的模可用常数倍的 $\exp(-\beta E_{g,\mathbf n}/2)$ 控制，因为 $\|\exp(-\eta H_e)\|$ 在这段区间内有共同上界。这个可求和的控制允许将 $\eta\downarrow0$ 与求迹交换；$\exp(-\eta H_e)$ 的强极限为恒等算符，遂恢复\eqnrefs{eq:thermal-correlation}。

若 $H_e=H_g$，热核复合律使积分直接等于 $Z_g$，故 $C(t)=1$。一般 $H_e$ 在 $t=0$ 时其核趋于恒等算符，仍给 $C(0)=1$；固定 $\eta>0$ 时的两个实热核和积分均为正，这就确定了连续平方根的起始支。
\end{proof}

\begin{proposition}[等频率位移模的有限温度边带]
取单模 $H_e=E_{eg}+D(\lambda)H_gD^\dagger(\lambda)$，$S=\lambda^2$，$\bar n=(\ee^{\beta\hbar\omega}-1)^{-1}$，则
\[
 C(t)=\ee^{-\ii E_{eg}t/\hbar}
 \exp\!\left[-S\{(\bar n+1)(1-\ee^{-\ii\omega t})
                     +\bar n(1-\ee^{\ii\omega t})\}\right].
\]
在 $\bar n>0$ 时，$C(t)=\ee^{-\ii E_{eg}t/\hbar}\sum_{k\in\mathbb Z}w_k\ee^{-\ii k\omega t}$，其中
\[
 w_k=\ee^{-S(2\bar n+1)}
 \left(\frac{\bar n+1}{\bar n}\right)^{k/2}
 I_{|k|}\!\left(2S\sqrt{\bar n(\bar n+1)}\right).
\]
$I_k$ 为第一类修正 Bessel 函数。权重满足 $\sum_kw_k=1$ 与 $w_{-k}=\ee^{-\beta\hbar\omega k}w_k$。只有 $\bar n\to0$ 时才退化为 $k\ge0$ 的 Poisson 分布。
\end{proposition}
\begin{proof}
利用 $\ee^{\ii H_gt/\hbar}D(\lambda)\ee^{-\ii H_gt/\hbar}=D(\lambda\ee^{\ii\omega t})$，把热关联中的两个演化合成两个位移算符。位移乘法关系
$D(\alpha)D(\gamma)=\exp[(\alpha\gamma^*-\alpha^*\gamma)/2]D(\alpha+\gamma)$
给出相位 $-\ii S\sin\omega t$。热态特征函数
$\Tr[\rho_gD(\alpha)]=\exp[-(\bar n+1/2)|\alpha|^2]$
则给出实部 $-S(2\bar n+1)(1-\cos\omega t)$，合并后得到 $C(t)$。

再将两个指数 $\exp[S(\bar n+1)\ee^{-\ii\omega t}]$ 和 $\exp[S\bar n\ee^{\ii\omega t}]$ 展开。保留净振动量子数 $k=p-q$ 的项，得到正项和
\[
 w_k=\ee^{-S(2\bar n+1)}
 \sum_{\substack{p,q\ge0\\p-q=k}}
 \frac{[S(\bar n+1)]^p(S\bar n)^q}{p!q!}.
\]
将此和与 $I_k$ 的幂级数比较，便得到 Bessel 形式。交换 $p,q$ 给出正负边带的权重比；令 $t=0$ 给出归一化。$\bar n\to0$ 时只剩 $q=0$，恢复\eqnrefs{eq:fc-poisson}。
\end{proof}

\begin{workedexample}[同一位移模的热边带与零温边带]
取 $S=1$、$\bar n=1$，求中心线及相邻两条边带的权重，并与零温结果比较。

\textbf{解。}此时 $\beta\hbar\omega=\log2$，Bessel 自变量为 $2\sqrt2$。代入上式得
\[
 w_0=\ee^{-3}I_0(2\sqrt2)\simeq0.2117,\qquad
 w_1=\ee^{-3}\sqrt2 I_1(2\sqrt2)\simeq0.2385,
 \qquad w_{-1}=w_1/2\simeq0.1192.
\]
其余权重分配给 $|k|\ge2$ 的边带；负边带来自热态中已有的振动量子。零温时 $w_0=w_1=\ee^{-1}\simeq0.3679$，所有 $k<0$ 权重为零。因此位移相同而温度不同，既改变包络宽度，也改变允许出现的频率方向。
\end{workedexample}
